% =====================================================================
\section{Record of the procedure}\label{app:history}
% =====================================================================

Four stages are kept apart. The \emph{decision cutoff} of 31 May 2025 applies to the historical records used as inputs, with the approximations of Section~\ref{sec:data}. At the \emph{procedural freeze} on 17 November 2026, the plan, the alternatives, all candidates considered and the evaluation measures were saved with file fingerprints before the evaluation code opened any 2025/26 record, and the \emph{evaluation} followed the same day. A \emph{post-freeze correction} on 18 November 2026, made after the outcomes had been examined, fixed errors in how some rules had been applied; no rule, weight or limit changed, and all 56 corrected runs returned the same transactions, objective, spending and wages. The report uses the corrected runs. The \emph{additional runs} of 21 and 22 November 2026 answer questions raised in review. They are later analyses, not part of the freeze: repeating the base case reproduced the original recommendation, and the additional scenarios examine how the plan changes under different assumptions. The freeze documents the procedure; it does not make the study a prospective experiment or its information set free of later knowledge.

{\setstretch{1}
\begin{longtable}{@{}>{\raggedright\arraybackslash}p{0.17\linewidth} >{\raggedright\arraybackslash}p{0.79\linewidth}@{}}
\caption{When the rules were set, what the correction after the freeze changed and which later analyses were added. The project's dated decision log holds the full entries.}\label{tab:history}\\
\toprule
Date (2026) & Stage and what was decided or changed \\
\midrule
\endfirsthead
\toprule
Date (2026) & Stage and what was decided or changed \\
\midrule
\endhead
\bottomrule
\endlastfoot
16 November & \emph{Before any run with real players:} wage allowance, cover bounds, role and depth weights, fee multipliers 1.1 and 1.3, sale receipt 0.9, score components and weights, wage bands. \\
16--17 November & \emph{After provisional plans had been seen:} clubs that do not sell; rival, regular-starter and wealthy-club mark-ups; one blocked former United player; the free-agent and two-year rules; fee ceiling lowered from \pounds 100m to \pounds 75m. These are feasibility and pricing assumptions. One input correction: wide midfielders are wing-backs only in a back three or five (45 pool players' roles changed). \\
17 November & \emph{Freeze and evaluation:} plan, alternatives, the 281 candidates then considered and the evaluation measures saved with file fingerprints before the evaluation code ran. \\
18 November & \emph{Post-freeze correction:} clubs identified by number with reviewed youth-team names (17 candidates changed two-year status; eligible 314 to 329, base-run 281 to 296); undated academy players need a dated record before June 2023; repeatable inputs; wages fixed in fee-only and club-rule runs; evaluation by club (Hincapié 0.550 to 0.524, Evann Guessand 0.268 to 0.172, Senne Lammens 0.947 to 0.842). All 56 runs unchanged; original files verified. \\
21 November & \emph{Additional runs:} more wage allowances, seven arrivals, 16 and 18 non-home-grown places, no depth reward, forced sales of Mount and Højlund, a protected wage limit with a 10\% deviation (an assumption made for this analysis) and the linear relaxation (Table~\ref{tab:extraruns}). Repeating the base case reproduced the original plan; each scenario changes only its stated limit, weight or restriction, and the base specification is unchanged. \\
22 November & \emph{Later analyses, not part of the freeze:} a scenario that rules out both Hincapié and Éderson and decides the whole plan again, and an audit that counts the reference players with league appearances for more than one club in a season (no score recalculated). Both were added to the folder of the 21 November runs, which were repeated from the committed code with identical results. \\
\end{longtable}}

\FloatBarrier
